\chapter*{CONCLUSION}
\addcontentsline{toc}{chapter}{Conclusion}
\markboth{CONCLUSION}{}

Ce stage de six mois au sein du pôle Advanced Data \& Climate d'Allianz France a permis de mener à bien la première migration opérationnelle d'un datamart SAS vers une architecture Python/PySpark en contexte industriel. Conçue comme un prototype pilote du programme SAS EXIST, cette mission dépasse le simple exercice technique : elle constitue une démonstration de faisabilité, une méthodologie documentée, et un ensemble de livrables réutilisables immédiatement mobilisables pour les migrations à venir.

Sur le plan technique, l'objectif central, garantir la parité fonctionnelle avec le système SAS d'origine, a été atteint pour les pipelines PTF Mouvements et Capitaux. Les comparaisons réalisées sur la vision de référence (décembre 2025) font apparaître un alignement strict : mêmes volumes de lignes, mêmes distributions par pôle et par produit, même structure de colonnes. La démarche itérative adoptée --- cinq phases ordonnées, validation continue, résolution par l'investigation et non par contournement --- a permis d'identifier et de corriger des divergences subtiles qui n'auraient pas résisté à un simple test de totalisation. La migration ne s'est pas limitée à transcrire du code SAS en Python : elle l'a réarchitecturé selon les principes de l'architecture médaillon (Bronze/Silver/Gold), en externalisant la configuration, en documentant chaque règle de gestion, et en produisant un code modulaire conçu pour être étendu.

Sur le plan méthodologique, ce stage a mis en lumière les conditions nécessaires à une telle migration. L'absence de documentation initiale du code SAS a constitué le premier obstacle : la reconstitution des règles métier à partir du seul code source a représenté une part significative de l'investissement total. Ce constat plaide pour que les futures migrations bénéficient dès leur lancement de la documentation fonctionnelle produite dans ce prototype. De même, la disponibilité des données de test en environnement Azure s'est révélée un facteur critique : les protocoles de sécurité stricts d'Allianz ont induit des délais non anticipés dans la phase de validation, soulignant la nécessité d'intégrer cet aspect dans le planning des projets futurs.

Pour Allianz France, les bénéfices attendus de cette migration s'articulent autour de trois axes. La réduction des coûts opérationnels liés aux licences SAS représente l'argument économique immédiat. L'amélioration des performances grâce au calcul distribué sur Azure Databricks constitue un levier technique structurel. Enfin, et c'est peut-être le bénéfice le plus durable, la constitution d'une documentation complète du datamart Construction, inexistante avant ce projet, capitalise une connaissance métier qui était jusqu'alors uniquement encodée dans un code source de quinze mille lignes non maintenu.

Les perspectives ouvertes par ce prototype sont directes. Le programme SAS EXIST compte quatre autres marchés à migrer (Auto, Habitation, Santé, Épargne). Le framework développé --- modules de lecture génériques, infrastructure de logging, configuration externalisée, méthodologie de validation --- est conçu pour être réutilisé sans modification majeure. La prochaine étape naturelle est la finalisation de la validation de parité sur dix visions historiques, puis le lancement de la migration du deuxième datamart en capitalisant sur les enseignements de ce premier cycle.
